\documentclass[a4paper,12pt]{article}
\usepackage[T2A]{fontenc}
\usepackage[utf8]{inputenc}
\usepackage[russian]{babel}
\usepackage{amsmath,amssymb,mathtools}
\usepackage{PTSans}
\renewcommand{\familydefault}{\sfdefault}
\usepackage{icomma}
\usepackage[a4paper,margin=18mm,headsep=7mm,footskip=12mm]{geometry}
\usepackage{microtype}
\usepackage{xcolor}
\usepackage{tikz}
\usetikzlibrary{calc,arrows.meta,positioning,shapes.geometric}
\usepackage{tkz-euclide}
\usepackage{tabularx,booktabs,longtable}
\usepackage{enumitem}
\usepackage{fancyhdr}
\usepackage{setspace}
\usepackage{hyperref}
\usepackage{parskip}
\usepackage{needspace}

\definecolor{accent}{HTML}{1F6F78}
\definecolor{darkaccent}{HTML}{174C52}
\definecolor{textgray}{HTML}{2B3033}
\definecolor{softgray}{HTML}{6B7478}
\definecolor{rulegray}{HTML}{CDD6D8}
\definecolor{softaccent}{HTML}{EAF3F4}
\definecolor{warn}{HTML}{8A4B2B}

\hypersetup{colorlinks=true,linkcolor=darkaccent,urlcolor=darkaccent}
\setlength{\parindent}{0pt}
\setlength{\parskip}{0.45em}
\onehalfspacing
\setlist[itemize]{leftmargin=1.25em,itemsep=0.25em,topsep=0.25em}
\setlist[enumerate]{leftmargin=1.55em,itemsep=0.35em,topsep=0.35em}
\renewcommand{\arraystretch}{1.25}

\pagestyle{fancy}
\fancyhf{}
\fancyhead[L]{\small\color{softgray}Геометрия · ОГЭ}
\fancyhead[R]{\small\color{softgray}София Кальней}
\fancyfoot[C]{\small\color{softgray}\thepage}
\renewcommand{\headrulewidth}{0pt}\setlength{\headheight}{14.5pt}

\newcommand{\sectionline}{\par\vspace{0.2em}\noindent\color{rulegray}\rule{\linewidth}{0.6pt}\par\color{textgray}\vspace{0.45em}}
\newcommand{\key}[1]{\textcolor{darkaccent}{\textbf{#1}}}
\newcommand{\formula}[1]{\[\boxed{#1}\]}
\newcommand{\lessondate}{3 октября 2026 г.}

\begin{document}
\color{textgray}

% ---------- TITLE ----------
\thispagestyle{empty}
\begin{center}
\vspace*{25mm}
{\Large\color{softgray}Чек-лист}\par
\vspace{7mm}
{\fontsize{26}{31}\selectfont\bfseries\color{darkaccent}Высоты, площади и\\[2mm]биссектрисы треугольника}\par
\vspace{10mm}
{\large Геометрия · подготовка к ОГЭ}\par
\vspace{18mm}
{\large София Кальней}\par
\vspace{4mm}
{\normalsize \lessondate}\par
\vfill
{\small\color{softgray}Лёвин Артём Александрович эксклюзивно для Софии Кальней}\par
\vspace{8mm}
\begin{tikzpicture}[x=1cm,y=1cm]
  \draw[accent,line width=1.2pt] (0,0) .. controls (3,0.9) and (7,-0.9) .. (10,0);
  \draw[rulegray,line width=0.7pt] (1,-0.35) .. controls (4,0.3) and (6,-0.3) .. (9,-0.35);
\end{tikzpicture}
\end{center}
\clearpage

% ---------- CORE ----------
\section*{1. Главная идея занятия}
\sectionline
В задачах на высоты и биссектрисы удобно связывать один и тот же треугольник с несколькими известными формулами. Центральные инструменты занятия:
\begin{itemize}
  \item площадь треугольника и формула Герона;
  \item два прямоугольных треугольника после проведения высоты;
  \item площадь через две стороны и синус угла между ними;
  \item формулы длины биссектрисы.
\end{itemize}

\key{Практический ориентир.} Если в задаче появляется высота, сразу ищите способ выразить площадь или используйте прямоугольные треугольники. Если появляется биссектриса, отметьте равные половины угла и проверьте, какая формула использует известные данные.

\section*{2. Высота через площадь}
\sectionline
Пусть в $\triangle ABC$ высота $BH$ проведена к стороне $AC=a$.

\begin{center}
\begin{tikzpicture}[scale=0.78]
  \tkzDefPoint(0,0){A}
  \tkzDefPoint(8,0){C}
  \tkzDefPoint(2.5,4.2){B}
  \tkzDefPointBy[projection=onto A--C](B)\tkzGetPoint{H}
  \tkzDrawPolygon[line width=1pt,color=accent](A,B,C)
  \tkzDrawSegment[line width=1pt,color=darkaccent](B,H)
  \tkzMarkRightAngle[size=0.35,color=darkaccent](B,H,C)
  \tkzDrawPoints[color=accent,fill=accent,size=2.5](A,B,C,H)
  \tkzLabelPoints[below](A,H,C)
  \tkzLabelPoints[above](B)
  \tkzLabelSegment[below](A,C){$a$}
  \tkzLabelSegment[right](B,H){$h_a$}
\end{tikzpicture}
\end{center}

Если известна площадь $S$, то
\formula{S=\frac12 a h_a \quad\Longrightarrow\quad h_a=\frac{2S}{a}.}

Если известны три стороны $a,b,c$, сначала используйте формулу Герона:
\formula{p=\frac{a+b+c}{2},\qquad S=\sqrt{p(p-a)(p-b)(p-c)}.}

\Needspace{8\baselineskip}
\key{Пример занятия.} Для сторон $10$, $17$, $21$:
\[
p=24,\qquad S=\sqrt{24\cdot14\cdot7\cdot3}=84,
\]
поэтому высота к стороне $21$ равна
\[
h=\frac{2\cdot84}{21}=8.
\]

\section*{3. Та же высота через два прямоугольных треугольника}
\sectionline
Обозначим $AH=y$, тогда $HC=a-y$. Для той же высоты $BH=x$ получаем две теоремы Пифагора:
\[
\begin{cases}
x^2+y^2=AB^2,\\
x^2+(a-y)^2=BC^2.
\end{cases}
\]
Удобный ход --- \key{вычесть уравнения}. Тогда $x^2$ исчезает, находится $y$, после чего $x$ восстанавливается из любого уравнения.

Для треугольника со сторонами $10$, $17$, $21$:
\[
\begin{cases}
x^2+y^2=100,\\
x^2+(21-y)^2=289.
\end{cases}
\]
После вычитания получается $y=6$, затем $x^2=100-36=64$, значит $x=8$.

\key{Когда метод особенно полезен.} Когда высота делит фигуру на два прямоугольных треугольника с удобными сторонами или когда площадь вычислять громоздко.

% ---------- AREA + BISECTOR ----------
\section*{4. Площадь через две стороны и синус угла}
\sectionline
Если известны стороны $b$ и $c$ и угол $A$ между ними, то
\formula{S=\frac12 bc\sin A.}
Эта формула работает для любой пары сторон, если в формуле стоит \key{угол между выбранными сторонами}.

\textbf{Мини-пример.} Если $b=8$, $c=11$, $A=30^\circ$, то
\[
S=\frac12\cdot8\cdot11\cdot\sin30^\circ=22.
\]

\clearpage
\section*{5. Биссектриса: формула через половину угла}
\sectionline
Пусть $AK$ --- биссектриса угла $A$ в $\triangle ABC$.

\begin{center}
\begin{tikzpicture}[scale=0.82]
  \tkzDefPoint(0,0){B}
  \tkzDefPoint(8,0){C}
  \tkzDefPoint(3.0,4.4){A}
  \tkzDefLine[bisector](B,A,C)\tkzGetPoint{T}
  \tkzInterLL(A,T)(B,C)\tkzGetPoint{K}
  \tkzDrawPolygon[line width=1pt,color=accent](A,B,C)
  \tkzDrawSegment[line width=1pt,color=darkaccent](A,K)
  \tkzMarkAngle[size=0.65,color=darkaccent](B,A,K)
  \tkzMarkAngle[size=0.85,color=darkaccent](K,A,C)
  \tkzDrawPoints[color=accent,fill=accent,size=2.5](A,B,C,K)
  \tkzLabelPoints[above](A)
  \tkzLabelPoints[below](B,K,C)
  \tkzLabelSegment[left](A,B){$c$}
  \tkzLabelSegment[right](A,C){$b$}
  \tkzLabelSegment[right](A,K){$l_a$}
\end{tikzpicture}
\end{center}

Пусть $\angle A=2\alpha$. Сравнение площади всего треугольника с суммой площадей $\triangle ABK$ и $\triangle ACK$ даёт:
\[
\frac12 bc\sin 2\alpha
=\frac12 l_a(c+b)\sin\alpha.
\]
Используя $\sin2\alpha=2\sin\alpha\cos\alpha$, получаем
\formula{l_a=\frac{2bc\cos\frac A2}{b+c}.}

\key{Смысл вывода.} Одна и та же площадь выражается двумя способами: сразу для большого треугольника и как сумма площадей двух треугольников, образованных биссектрисой.

\section*{6. Биссектриса: формула через отрезки противоположной стороны}
\sectionline
Если биссектриса $AK$ делит сторону $BC$ на отрезки $BK=m$ и $KC=n$, то
\formula{AK^2=AB\cdot AC-BK\cdot KC.}
То есть
\[
l_a^2=bc-mn.
\]
Формула особенно удобна, когда длины $BK$ и $KC$ уже известны или легко находятся по теореме о биссектрисе:
\formula{\frac{BK}{KC}=\frac{AB}{AC}.}

\section*{7. Что обязательно проверять}
\sectionline
\begin{itemize}
  \item В формуле $S=\frac12 bc\sin A$ угол $A$ должен лежать \key{между} сторонами $b$ и $c$.
  \item Высота --- положительная длина, поэтому из $h^2=q$ берём $h=\sqrt q$.
  \item После проведения высоты подпишите оба отрезка основания: например, $AH=y$, $HC=a-y$.
  \item В формуле биссектрисы $l_a^2=bc-mn$ произведение $mn$ относится к двум частям стороны, которую пересекает биссектриса.
  \item Перед вычислениями проверьте, что стороны образуют треугольник.
  \item Корни и дроби можно оставлять в точном виде, если условие не требует десятичного приближения.
\end{itemize}
\clearpage

% ---------- FLOWCHART ----------
\thispagestyle{fancy}
\begin{center}
{\Large\bfseries\color{darkaccent}Алгоритм: как искать высоту треугольника}
\end{center}
\vspace{8mm}

\begin{center}
\begin{tikzpicture}[
  startstop/.style={ellipse,draw=accent,line width=1pt,minimum width=42mm,minimum height=10mm,align=center,fill=softaccent},
  process/.style={rectangle,rounded corners=2mm,draw=accent,line width=1pt,text width=40mm,minimum height=13mm,align=center},
  decision/.style={diamond,aspect=2.25,draw=accent,line width=1pt,text width=28mm,align=center,inner sep=1.2mm},
  arrow/.style={-{Stealth[length=3mm]},line width=0.9pt,color=darkaccent},
  lab/.style={font=\small,fill=white,inner sep=1pt}
]
\node[startstop] (start) at (0,0) {Нужно найти высоту $h_a$};
\node[decision] (three) at (0,-2.3) {Известны три стороны?};
\node[process] (heron) at (-4.0,-5.0) {Найдите площадь по формуле Герона};
\node[decision] (angle) at (3.5,-5.0) {Известны две стороны и угол между ними?};
\node[process] (fromarea) at (-4.0,-8.0) {Используйте $h_a=\dfrac{2S}{a}$};
\node[process] (sine) at (1.0,-8.3) {Найдите $S=\dfrac12 bc\sin A$};
\node[process] (right) at (6.0,-8.3) {Проведите высоту и составьте две теоремы Пифагора};
\node[startstop] (finish) at (1.0,-11.5) {Проверьте длину и запишите ответ};

\draw[arrow] (start) -- (three);
\draw[arrow] (three) -- node[lab,above left]{да} (heron);
\draw[arrow] (three) -- node[lab,above right]{нет} (angle);
\draw[arrow] (heron) -- (fromarea);
\draw[arrow] (angle) -- node[lab,above left]{да} (sine);
\draw[arrow] (angle) -- node[lab,above right]{нет} (right);
\draw[arrow] (sine) -- (finish);
\draw[arrow] (right) -- (finish);
\draw[arrow] (fromarea) -- (finish);
\end{tikzpicture}
\end{center}
\clearpage

% ---------- TRAINING ----------
\section*{Тренировочный блок · Путь А}
\sectionline
\textbf{10 задач.} Решайте по порядку: от прямого применения формул к сочетанию нескольких идей.

\begin{enumerate}
  \item В треугольнике стороны равны $6$, $8$ и $10$. Найдите высоту, проведённую к стороне $10$.

  \item В треугольнике стороны равны $13$, $14$ и $15$. Найдите высоту, проведённую к стороне $14$, используя формулу Герона.

  \item В треугольнике $ABC$ даны $AB=15$, $BC=20$, $AC=25$. Высота $BH$ проведена к $AC$. Решите задачу через два прямоугольных треугольника: найдите $AH$ и $BH$.

  \item Две стороны треугольника равны $8$ и $11$, угол между ними равен $30^\circ$. Найдите площадь треугольника.

  \item Площадь треугольника равна $30$. Одна сторона равна $12$, а высота проведена к этой стороне. Найдите высоту.

  \item В треугольнике две стороны, образующие угол $60^\circ$, равны $6$ и $8$. Найдите длину биссектрисы этого угла по формуле через косинус половины угла.

  \item В $\triangle ABC$ биссектриса $AK$ проведена к стороне $BC$. Известно: $AB=10$, $AC=15$, $BC=20$. Найдите $BK$, $KC$ и $AK$.

  \item Стороны треугольника равны $7$, $8$ и $9$. Найдите высоту, проведённую к стороне $9$. Ответ запишите в точном виде.

  \item Две стороны треугольника равны $8$ и $10$, угол между ними равен $120^\circ$. Найдите площадь треугольника.

  \item В треугольнике со сторонами $13$, $14$, $15$ к стороне $14$ проведены высота и биссектриса из противоположной вершины. Найдите длины обеих линий.
\end{enumerate}
\clearpage

% ---------- ANSWERS ----------
\section*{Ответы}
\sectionline
\begin{longtable}{@{}>{\bfseries}p{10mm}p{0.80\linewidth}@{}}
\toprule
№ & Ответ \\
\midrule
\endhead
1 & $\dfrac{24}{5}=4{,}8$. \\
2 & $12$. \\
3 & $AH=9$, $BH=12$. \\
4 & $22$. \\
5 & $5$. \\
6 & $\dfrac{24\sqrt3}{7}$. \\
7 & $BK=8$, $KC=12$, $AK=3\sqrt6$. \\
8 & $\dfrac{8\sqrt5}{3}$. \\
9 & $20\sqrt3$. \\
10 & высота $12$, биссектриса $\dfrac{3\sqrt{65}}{2}$. \\
\bottomrule
\end{longtable}

\vspace{8mm}
\textcolor{softgray}{\small Самопроверка: если ответ на высоту превышает обе прилежащие стороны, вернитесь к вычислению площади или к теореме Пифагора.}

\end{document}
